\chapter{Designing and Implementing SpeechPrint}

\section{Initial Prototype}

\section{Implementing SpeechPrint-V2}

\subsection{Selecting Musical Features for SpeechPrint-V2}

%\subsection{Designing SpeechPrint}

%\section{SpeechPrint System Architecture}


%\chapter{SpeechPrint Evaluation Study}

%\chapter{Methods}

\chapter{Methods}
\label{methods}

\section{Pitch Extraction and Octave Error Correction}
\label{sec:pitch_extraction}

A central challenge in automated prosody labelling is the reliability of the
underlying fundamental frequency (F0) measurements. SpeechPrint originally used
Praat's default autocorrelation pitch tracker (the SCC method), which is
susceptible to \emph{octave errors}: frames where the tracker identifies a
harmonic partial rather than the true F0. An octave error is not a small
perturbation --- it corresponds to a shift of exactly 12 semitones, which is
typically the full pitch range of a speaker. When the F0 onset or offset of a
syllable nucleus is contaminated by such an error, the intra-syllable pitch
movement $\Delta_{F_0}$ can be completely wrong. For example, if the true F0 at
syllable onset is 214\,Hz and Praat mistakenly reports 107\,Hz (the second
sub-harmonic, 12\,ST below), a genuinely rising syllable will appear to the
labeller as level or even falling. As demonstrated empirically below, 238 such
octave errors were found in the first 30 seconds of the English test recording
alone --- directly explaining many of the erroneous prosody labels previously
observed.

\subsection{Comparative Evaluation: Praat, Librosa pYIN, and CREPE}
\label{sec:tracker_comparison}

To address this, three pitch trackers were evaluated and compared:

\begin{enumerate}

\item \textbf{Praat SCC (raw):} The default Praat pitch tracker, using the
      cross-correlation method with a pitch floor of 75\,Hz and ceiling of
      600\,Hz. This is the out-of-the-box Parselmouth call (\texttt{to\_pitch}).
      It is susceptible to octave jumps, particularly at voicing boundaries,
      because the autocorrelation function has local maxima at integer multiples
      of the period.

\item \textbf{Praat AC + Xu (1999) correction:} An improved Praat call
      (\texttt{to\_pitch\_ac}) with \texttt{octave\_jump\_cost=0.5} (double the
      default penalty), which discourages the tracker from jumping an octave
      between adjacent frames. Residual outliers are then removed using the
      median-based spike-removal algorithm of Xu~(1999) \cite{xu1999}: the median
      F0 across all voiced frames is computed; any frame whose F0 deviates more
      than 12 semitones from that median is replaced with \texttt{None}. The
      surviving values are then smoothed with one pass of triangular weighting
      (weights 1:2:1 over three consecutive voiced frames), attenuating
      frame-to-frame noise while preserving the overall pitch trajectory.

\item \textbf{Librosa pYIN (probabilistic YIN):} A state-of-the-art monophonic
      F0 tracker based on the probabilistic extension of the YIN algorithm
      \cite{mauch2014pyin}. Unlike Praat, pYIN returns a voicing probability for
      each frame, allowing unvoiced frames to be cleanly masked rather than
      forced to zero. It is less susceptible to octave errors because it jointly
      models F0 continuity and voicing likelihood. The implementation from
      librosa 0.11 \cite{mcfee2015librosa} was used with a frame length of 2048
      samples and a hop of 256 samples at 16\,kHz, giving a temporal resolution
      of approximately 16\,ms.

\item \textbf{CREPE:} The Convolutional REpresentation for Pitch Estimation
      \cite{kim2018crepe} is a data-driven pitch tracker published at ICASSP 2018
      by the Music and Audio Research Laboratory (MARL) at NYU Steinhardt. Unlike
      signal-processing methods such as Praat and pYIN, CREPE does not model the
      vocal source as a harmonic series and does not compute an autocorrelation or
      difference function. Instead, it feeds the raw time-domain waveform (a
      1024-sample window centred on each frame) directly into a six-layer deep
      convolutional neural network, which produces a 360-bin probability
      distribution over log-frequency bins spanning 32.7\,Hz to 1975.5\,Hz. The
      predicted F0 is the expected value of this distribution; the maximum bin
      confidence serves as a voicing probability, replacing the need for a
      separate voiced/unvoiced detector. Because CREPE learns its pitch model
      from a large and diverse training corpus rather than deriving it from
      acoustic first principles, it is not fooled by irregular phonation patterns
      (creaky voice, breathy voice, laryngealisation) that can corrupt the
      autocorrelation peak used by Praat and pYIN. The PyTorch port
      \texttt{torchcrepe} \cite{morrisonTorchcrepe} was used with the
      \texttt{full} model capacity, a 10\,ms hop length (giving a temporal
      resolution of 16,269 frames for the 162.7\,s DoReCo recording), batched
      inference (\texttt{batch\_size=128}) to stay within CPU memory limits, and a
      periodicity confidence threshold of 0.5 for the voiced/unvoiced decision.
      CREPE requires resampling the audio to 16\,kHz before inference; this step
      is performed internally by \texttt{torchcrepe}.

\end{enumerate}

\subsection{Speaker-adaptive pitch range}

A prerequisite for reliable pitch extraction is a pitch search range matched to
the speaker. The default Praat range (75--600\,Hz) covers both bass and soprano
voice types simultaneously, which increases the risk of octave errors in
narrower-range recordings. SpeechPrint v3 therefore auto-detects the speaker's
pitch range from a 20-second scan of the recording using pYIN at loose bounds
(65--500\,Hz), computes the median of the voiced frames, and sets the working
floor and ceiling to $\text{median} \times 0.5$ and $\text{median} \times 2.5$
respectively. This restricts the search space to a plausible single-speaker range
without manual tuning.

\subsection{Final pipeline}

Two pipeline configurations are used depending on the data source.

For the \textbf{English and German recordings} (ASR-based pipeline), Librosa
pYIN is the primary pitch tracker. pYIN requires no model download, runs at
approximately 3--5$\times$ real-time on CPU, and its voiced/unvoiced probability
mask cleanly suppresses unvoiced frames without further intervention. The full
pitch track is extracted once per audio file and stored in memory; subsequent
syllable nucleus measurements look up the pre-computed values at the relevant
time points. Xu~(1999) spike removal is applied as a safety net for any residual
octave artefacts.

For the \textbf{DoReCo Daakie endangered-language recording} (human-annotation
pipeline), CREPE is used as the primary pitch tracker, with the rationale
developed in this section: CREPE produces the most accurate F0 estimates on
challenging material, it does not rely on harmonic assumptions that may be
violated by the phonation characteristics of under-resourced languages, and the
computational cost is acceptable when processing a single archival recording
offline. Xu~(1999) spike removal is applied after CREPE as a safety net, and the
relative prosody labelling procedure described in
Section~\ref{sec:prosody_label_method} is applied identically to both
pipelines.

% ADDED HERE — implementation, methodology, evaluation design, demo design, amplitude, relative vs absolute
% (Lonce feedback: write what you actually did; more text beyond literature review; demo design with ASCII; methodology)

\section{System Architecture and Implementation}
\label{sec:implementation}

SpeechPrint is implemented as a Python pipeline organized into nine sequential
stages. The pipeline is invoked from the command line with a WAV file and a
language code; it produces a Praat-format TextGrid, a set of CSV tables, and a
JSON manifest, all written to a per-recording output directory. The nine stages
are as follows.

\paragraph{Stage 1: Load audio.}
The WAV file is read using the Python \texttt{wave} module to extract sample
rate, channel count, and total duration. These values are stored in the manifest
and used to validate downstream timing assertions.

\paragraph{Stage 2: Transcribe.}
WhisperX with the \texttt{large-v3} model is used to obtain a text transcript
and coarse segment timing. The model is loaded from a local cache to avoid
repeated downloads. Voice activity detection parameters
(\texttt{vad\_onset=0.3}, \texttt{vad\_offset=0.2}) are lowered relative to the
defaults in order to detect short initial syllables that would otherwise be
missed. If WhisperX is unavailable, the system falls back to plain Whisper, and
if that is also unavailable, to a filename-derived fallback transcript.

\paragraph{Stage 3: Prepare transcript.}
The raw transcript is cleaned of punctuation, normalized to lowercase, and split
into word tokens. Silent intervals longer than 40\,ms between consecutive word
boundaries are labelled \texttt{<sil N.NNs>} in the words tier.

\paragraph{Stage 4: Forced alignment.}
For English, the Montreal Forced Aligner (MFA) is used when available, providing
phone-level timing. For other languages, WhisperX's CTC alignment step is used,
which gives word-level timestamps. If neither is available, words are distributed
proportionally within each Whisper segment.

\paragraph{Stage 5: Pitch and intensity extraction.}
A single full-file pitch track is extracted using Librosa pYIN (see
Section~\ref{sec:pitch_extraction}) and stored in memory. Intensity is extracted
using Parselmouth's \texttt{to\_intensity} method with a 40\,ms window. Both
tracks are computed once and subsequently queried by time index for each syllable
measurement, avoiding redundant file I/O.

\paragraph{Stage 6: Phonemization.}
Word strings are phonemized using \texttt{espeak-ng} via the Python phonemizer
library \cite{barnard2021phonemizer}, which produces IPA phone sequences. The
language code is mapped to an espeak backend code (e.g.\ \texttt{en-us} for
English, \texttt{de} for German). Phones are then grouped into syllables using a
maximum-onset heuristic: vowel nuclei are identified by membership in a fixed set
of IPA vowel characters, and intervocalic consonant clusters are split so that as
many consonants as possible are assigned to the following syllable onset.

\paragraph{Stage 7: Symbolic prosody labelling.}
Prosody labels are assigned in two sub-steps. First, acoustic measurements are
made at each syllable's vowel nucleus: ten time-evenly-spaced F0 values are
sampled from the pre-computed pYIN track, and the mean intensity over the same
window is read from the Parselmouth intensity object. Second, the labelling
function (see Section~\ref{sec:prosody_label_method}) assigns a symbol to each
syllable based on the computed measurements and the measurements of its
neighbours.

\paragraph{Stage 8: Write TextGrid.}
A five-tier Praat TextGrid is written in plain text format. Empty intervals are
inserted automatically to ensure each tier covers the full duration of the
recording contiguously, as required by Praat's parser.

\paragraph{Stage 9: Export CSV and JSON.}
Per-word, per-syllable, and per-phoneme tables are exported as CSV files for
downstream statistical analysis. A recording-level prosody summary table and a
JSON manifest containing all extracted values are also written.

\section{Syllabification and Nucleus Identification}
\label{sec:syllabification}

Syllable boundaries within each word are determined by the positions of vowel
nuclei in the IPA phone sequence. The vowel detection function checks membership
in a fixed set of IPA vowel symbols that covers cardinal monophthongs,
IPA-specific vowel characters (schwa, open-mid, near-close, etc.), long-vowel
digraphs, and common English diphthong sequences. Consonantal glides (/w/, /j/,
/l/, /r/, etc.) are explicitly excluded even when they contain vowel-like
acoustic properties.

For each pair of consecutive vowels separated by one or more consonants, the
consonant cluster is split at the midpoint, following the maximum-onset principle
of syllable theory \cite{kahn1976syllable}: the second half of the cluster forms
the onset of the following syllable, and the first half forms the coda of the
preceding syllable. Where MFA phone-level timing is available, the actual onset
and offset times of the vowel phone are used directly for the nucleus window.
Otherwise, the nucleus window is estimated proportionally within the syllable
span.

\section{Prosody Labelling Algorithm}
\label{sec:prosody_label_method}

The symbolic prosody labels are assigned using the following procedure, applied
to each syllable in sequence.

\paragraph{Pitch movement.}
The intra-syllable pitch movement $\Delta_{F_0}$ is computed as:
\begin{equation}
  \Delta_{F_0} = 12 \log_2\!\left(\frac{F_0^{\mathrm{offset}}}{F_0^{\mathrm{onset}}}\right)
\end{equation}
where $F_0^{\mathrm{onset}}$ is the first voiced pYIN sample in the nucleus
window and $F_0^{\mathrm{offset}}$ is the last voiced sample. Pitch movement is
considered \emph{rising} if $\Delta_{F_0} > \theta_{\mathrm{weak}}$ and
\emph{falling} if $\Delta_{F_0} < -\theta_{\mathrm{weak}}$.

\paragraph{Adaptive thresholds.}
The weak-rise threshold $\theta_{\mathrm{weak}}$ is computed from the standard
deviation of all intra-syllable pitch movements in the recording:
\begin{equation}
  \theta_{\mathrm{weak}} = \max\!\left(\theta_{\mathrm{floor}},\;
    \alpha \cdot \mathrm{std}(\{\Delta_{F_0,i}\})\right)
\end{equation}
with a floor of $\theta_{\mathrm{floor}} = 0.5$\,ST and an adaptive factor of
$\alpha = 0.35$. A strong movement is defined as
$|\Delta_{F_0}| \geq 2.5 \cdot \theta_{\mathrm{weak}}$, and is labelled
\texttt{//} or \texttt{\textbackslash\textbackslash} rather than \texttt{/} or
\texttt{\textbackslash}. This design ensures that the thresholds scale with the
actual pitch range of the recording rather than applying a fixed semitone
criterion regardless of speaker.

\paragraph{Height relative to neighbours.}
In addition to the intra-syllable movement, the mean F0 of each syllable is
compared to the mean of its immediate neighbours. A syllable whose mean F0 is
more than 0.8\,ST above the neighbourhood mean is labelled as high
(\texttt{{\={}}}, U+203E overline); one more than 0.8\,ST below is labelled as
low (\texttt{\_}). The overline was chosen as a visual counterpart to the
underscore: both are single horizontal strokes, one above and one below the
text baseline, making the high/low contrast immediately readable in a Praat tier.

\paragraph{Relative amplitude and accent detection.}
The amplitude of each syllable is compared to the mean amplitude of its
neighbours in the same $\pm 2$ syllable window. If the syllable is more than
1.5\,dB louder than its neighbours and simultaneously higher in mean F0 (by more
than 1.0\,ST) or longer in duration (by a factor of at least 1.25), it receives
an accent marker \texttt{*}. This three-cue definition of prominence (louder,
higher, longer) is motivated by the multi-parametric nature of prosodic
prominence documented in the phonetic literature \cite{ladd2008intonational}.

\paragraph{Symbol assembly.}
The final symbol is assembled as \texttt{[*][height][direction]}, where each
component is present only if the corresponding condition is satisfied. For
example, a syllable that is prominent, high, and rising receives
\texttt{*{\={}}//}; one that is low and weakly falling receives
\texttt{\_\textbackslash}; one that is merely falling with no height contrast
receives \texttt{\textbackslash}; and one for which no voiced frames are
available receives \texttt{?}.

\section{Amplitude Analysis: Window Size and Relative Measurement}
\label{sec:amplitude}

A specific concern raised during evaluation was whether the amplitude analysis
window was appropriate. The original implementation measured mean intensity over
the entire syllable duration, which for syllables spanning 200--400\,ms could
smooth over the rapid amplitude rise that marks the onset of a prominent accent.
The supervisor noted: ``you'd want an RMS follower which would operate on a
shorter window of time and then you'd smooth it'' (Krifka, meeting June 2026).

In SpeechPrint v3, intensity is extracted using Parselmouth with a 40\,ms
analysis window (reduced from the Praat default of approximately 100\,ms) and
then averaged over the nucleus window. This provides a more responsive estimate
of the amplitude peak within the syllable rather than the mean across the whole
duration including quieter onset consonants.

The key change in the labelling logic is the shift from absolute to relative
amplitude. The original system compared amplitude against a fixed threshold in
decibels. The revised system computes the difference between the syllable's mean
intensity and the mean intensity of its neighbours within a $\pm 2$ syllable
window, which produces a speaker-normalised and recording-normalised quantity.
This addresses the concern that ``66 decibels doesn't sound outrageous but I
don't know in the context of what the other ones are'' (Krifka, meeting June 2026):
what matters for prominence is not the absolute level but the local contrast.

\section{Evaluation Design}
\label{sec:eval_design}

The evaluation of SpeechPrint's prosody output takes two forms.

\subsection{Automatic benchmark: GToBI comparison}

The five German sentences from the GToBI training corpus
(\url{http://www.gtobi.uni-koeln.de/}) were used as a controlled benchmark.
These sentences were selected because they have been annotated by expert
phoneticians according to GToBI conventions and represent a range of tone types:
L+H* (rising nuclear accent), H+L* (falling), H+!H* (downstepped high), and
L*+H (rising accent with low starred tone). The comparison procedure is:
\begin{enumerate}
\item Human word boundaries from the Wort tier are used as input, removing
      alignment error as a confound.
\item SpeechPrint's prosody labels are assigned using the procedure described
      in Section~\ref{sec:prosody_label_method}.
\item The GToBI tone symbol and its position within the word sequence are
      compared to the SpeechPrint label on the corresponding syllable.
\end{enumerate}

\subsection{Perceptual evaluation: questionnaire}

A short questionnaire will be administered to linguists who are not specialists
in phonetics. Participants are shown a Praat TextGrid display of a sentence
alongside an audio recording and asked to judge: (a) whether the prosody tier
symbols match their perception of the intonation, (b) whether the accent marker
appears on the correct word, and (c) how useful they would find such an
annotation in their own work. The questionnaire covers English minimal pairs and
selected German GToBI sentences.

\section{Demo Design: ASCII Prosody Symbols as Visual Representation}
\label{sec:demo_design}

A design principle of SpeechPrint is that the prosody representation should be
legible without specialist software. The chosen symbols are:

\begin{center}
\begin{tabular}{cl}
\hline
Symbol & Meaning \\
\hline
\texttt{/}  & Weakly rising pitch within the syllable \\
\texttt{//} & Strongly rising \\
\texttt{\textbackslash}  & Weakly falling pitch \\
\texttt{\textbackslash\textbackslash} & Strongly falling \\
\texttt{{\={}}} & High level (mean F0 above neighbouring syllables) \\
\texttt{\_} & Low level (mean F0 below neighbouring syllables) \\
\texttt{*}  & Prominent accent (louder, higher, or longer than neighbours) \\
\texttt{?}  & Unvoiced or insufficient F0 data \\
\hline
\end{tabular}
\end{center}

Symbols may be combined: for example, \texttt{*{\={}}//} denotes a prominently
accented syllable that is high in the speaker's range and strongly rising; \texttt{\_\textbackslash\textbackslash}
denotes a syllable that is low in the speaker's range and strongly falling.

These symbols were chosen for visual legibility without specialist software.
\texttt{/} and \texttt{\textbackslash} resemble the slopes of a pitch contour;
the overline \texttt{{\={}}} and underscore \texttt{\_} form a natural pair
marking height above and below the local mean; \texttt{*} is a conventional
prominence marker in linguistic notation. All symbols are available in plain
ASCII and render unambiguously in a monospaced font and in Praat tier columns.

The demo video for the evaluation questionnaire walks through the following
sequence: (1) a raw audio recording is shown in Praat; (2) SpeechPrint
processes it and adds five tiers; (3) the prosody tier is highlighted and the
symbols are explained; (4) two minimal-pair sentences differing only in which
word receives the accent are compared side by side, showing how the \texttt{*}
marker moves between tiers. The demo is kept to approximately five minutes, in
line with the supervisor's recommendation that longer demos cause participants
to lose focus before reaching the questionnaire items.

\section{Linguistic Rationale for Syllable-Level Annotation}
\label{sec:linguistic_rationale}

The choice to annotate prosody at the syllable level rather than the word or
phoneme level reflects several considerations from the linguistic literature.
Prosodic features such as pitch accents and boundary tones are canonically
associated with syllables in metrical and autosegmental phonology
\cite{ladd2008intonational, beckman2006tobi}. In languages with lexical stress,
such as English and German, the accented syllable within a word is the primary
site of F0 movements and amplitude peaks. Annotating at the word level would
obscure which part of a multi-syllable word carries the prosodic event; annotating
at the phoneme level would over-specify, since prosodic movements span the entire
nucleus and not just a single phone.

Vowel nuclei receive special attention because they are the primary carriers of
F0: the vocal folds can vibrate continuously through a vowel, whereas consonants
frequently produce brief interruptions or reductions in periodicity. This means
that a pitch measurement made during a fricative or stop consonant would often
return undefined values, whereas a measurement made during the vowel nucleus is
almost always reliable. The ten-point time-normalised sampling procedure adopted
here, in which F0 is sampled at equal intervals across the nucleus, follows the
algorithm described in Xu~(1999) \cite{xu1999} and allows comparison of pitch
contour shapes across syllables of different durations.

\section{Algorithmic Optimisations to the Prosody Labeller}
\label{sec:optimisations}

The baseline prosody labelling algorithm described in
Section~\ref{sec:prosody_label_method} was improved in four ways, motivated by
observable deficiencies in the output distribution rather than theoretical
assumptions. Each optimisation targets a specific, diagnosable failure mode.

\subsection{MAD-based adaptive threshold}

The baseline uses the standard deviation $\sigma$ of all intra-syllable pitch
movements to set the weak-rise threshold $\theta_{\mathrm{weak}} = \max(0.5,
0.35\sigma)$. Standard deviation is sensitive to outliers: a small number of
syllables with very fast pitch movements (e.g.\ a sharp accent) inflates
$\sigma$, which in turn inflates both thresholds and causes many genuinely
strong movements to be classified as weak.

The fix is to replace $\sigma$ with the \emph{normalised median absolute
deviation} (MAD):
\begin{equation}
  \hat{\sigma}_{\mathrm{MAD}} = 1.4826 \cdot \mathrm{median}\!\left(
    \{|\Delta_{F_0,i} - \tilde{\Delta}|\}\right)
\end{equation}
where $\tilde{\Delta}$ is the median pitch movement and the factor 1.4826
makes $\hat{\sigma}_{\mathrm{MAD}}$ a consistent estimator of $\sigma$ under a
Gaussian assumption \cite{rousseeuw1993alternatives}. MAD is the standard robust
alternative to $\sigma$ used in outlier detection. The empirical effect on the
Daakie recording: the fraction of syllables labelled as \emph{strongly} rising
or falling fell from 30.9\% to 21.0\%, moving from an implausibly high rate
toward the expected range of 15--20\%.

\subsection{Utterance boundary reset}

The baseline compares each syllable's mean F0 and amplitude to those of its
immediate neighbours without regard for utterance structure. This causes errors
at utterance boundaries: the final syllable of an utterance typically undergoes
\emph{final lowering} (a well-documented phonetic phenomenon whereby F0 drops
toward the end of an intonational phrase), making it systematically lower in F0
and amplitude than the first syllable of the following utterance, which resets
to a higher level. Comparing these two syllables as neighbours therefore produces
a spurious high-versus-low contrast.

The fix is straightforward: the end time of each utterance is extracted from the
\texttt{tx@TA} tier and used as a \emph{break point}. Syllables on opposite
sides of a break are never used as each other's neighbours in the height or
amplitude comparisons. Only within-utterance neighbour pairs are considered.

\subsection{Pitch declination removal}

Even within a single utterance, F0 exhibits a systematic downward trend over
time known as \emph{declination}: the baseline F0 at the start of an utterance
is higher than at the end, independent of the prosodic structure
\cite{pierrehumbert1980phonology}. Without correction, this trend inflates the
count of syllables labelled as falling simply because their F0 is lower than
the preceding syllable---a consequence of the global trend, not of a genuine
falling movement.

The fix is to fit a linear regression to the sequence of syllable-mean F0
values within each utterance as a function of time:
\begin{equation}
  \hat{f}(t) = a \cdot t + b
\end{equation}
and to subtract the slope component before the height comparison:
\begin{equation}
  f_0^{\mathrm{det}}(i) = f_0^{\mathrm{mean}}(i) - a \cdot (t_i - \bar{t})
\end{equation}
where $\bar{t}$ is the mean time within the utterance. Subtracting only the
slope (not the intercept) preserves the absolute F0 level for the accent
detector, while removing the temporal drift. The empirical effect on the Daakie
recording: the falls-to-rises ratio improved from 0.61 (26.7\% falls,
16.3\% rises) to 0.72 (18.8\% falls, 13.5\% rises), indicating a more balanced
and plausible prosodic label distribution.

\subsection{Nucleus edge trimming}

The baseline samples F0 at ten equally spaced points across the full vowel
nucleus window $[t_{\mathrm{v0}}, t_{\mathrm{v1}}]$. The first and last
sampling points are therefore positioned at the very onset and offset of the
vowel, where the vocal tract is still transitioning from or to the adjacent
consonant. This transition region typically produces lower F0 values (due to
incomplete glottal adduction) or missing voicing, which contaminates the onset
and offset estimates used to compute the intra-syllable pitch movement.

The fix is to trim 15\% from each end of the nucleus window before sampling:
\begin{equation}
  t_{\mathrm{start}} = t_{\mathrm{v0}} + 0.15 \cdot d, \quad
  t_{\mathrm{end}}   = t_{\mathrm{v1}} - 0.15 \cdot d
\end{equation}
where $d = t_{\mathrm{v1}} - t_{\mathrm{v0}}$ is the nucleus duration. The
trimmed window is used only if it is at least 10\,ms long; otherwise the full
window is retained. The empirical effect is a slight increase in the unknown
rate (from 0.9\% to 1.5\%), which is expected: more conservative sampling means
a short or partially devoiced nucleus is more likely to be flagged as
unvoiced---a correct outcome rather than a spurious one.

\subsection{Octave recovery in spike removal}

The baseline Xu~(1999) implementation sets any F0 frame deviating more than
12\,ST from the utterance median to \texttt{None}, discarding it. However, many
such frames are not random noise but genuine octave errors: the tracker found
the second harmonic at $f_0 / 2$ or a doubled pitch at $2f_0$ instead of the
true fundamental. The octave alternative is the correct value and should be
recovered rather than discarded.

The fix: before setting a flagged frame to \texttt{None}, the two octave
alternatives $f \times 2$ and $f \div 2$ are tested against the utterance
median. If either is within 12\,ST of the median, the closest one replaces the
original frame. Only if both alternatives also fail the threshold is the frame
discarded. This recovers data points that would otherwise inflate the unknown
rate.

\section{Application to DoReCo Daakie: Endangered-Language Pipeline}
\label{sec:doreco_pipeline}

The DoReCo corpus provides a 162.7-second archival recording of Daakie (Port
Vato, Vanuatu) with detailed human-annotated TextGrid tiers \cite{seifart2022doreco}.
Because no forced aligner or ASR model exists for Daakie, the standard
SpeechPrint word-alignment and phonemization stages (Stages 2--4 and 6) are
replaced entirely by the existing human annotations, and only the acoustic
measurement and prosody-labelling stages are applied automatically.

\subsection{Tier sources}

The following tiers are taken directly from the DoReCo TextGrid without
modification:

\begin{description}
  \item[\texttt{sentence}] Utterance-level Daakie orthographic transcription
        (\texttt{tx@TA} tier).
  \item[\texttt{translation}] Utterance-level English free translation
        (\texttt{ft@TA} tier).
  \item[\texttt{words}] Daakie word intervals (\texttt{wd@TA} tier). Silent
        intervals longer than 40\,ms are labelled \texttt{<sil N.NNs>}.
  \item[\texttt{gloss}] English morpheme glosses (\texttt{gl@TA} tier),
        aligned to the same word intervals.
  \item[\texttt{phones}] SAMPA-based phoneme intervals (\texttt{ph@TA} tier),
        as annotated by the DoReCo transcribers. The Daakie phone inventory
        includes length-marked vowels (\texttt{a:}, \texttt{e:}, etc.),
        labialized consonants (\texttt{b\_w}, \texttt{m\_w}), and the
        near-open front vowel \texttt{\{}.
\end{description}

\subsection{Syllabification from human-annotated phones}

Because no Daakie syllabifier exists, syllable boundaries are derived from the
\texttt{ph@TA} phone sequence using the same maximum-onset heuristic applied to
IPA in the English pipeline. The Daakie vowel set is defined by membership in a
fixed inventory established from the full phone tier. Each word's phones are
grouped into syllables by identifying vowel nuclei and splitting intervocalic
consonant clusters at the midpoint. The nucleus window for F0 measurement is the
span of vowel phones within the syllable, as determined by the human phone-level
timing.

\subsection{F0 extraction: CREPE at full resolution}

CREPE (\texttt{torchcrepe}, \texttt{full} model, 10\,ms hop) is used as the
F0 tracker for the DoReCo pipeline rather than pYIN, for the reasons established
in Section~\ref{sec:tracker_comparison}: CREPE makes no assumptions about the
harmonic structure of the voice source, which is particularly important for an
under-resourced language where the phonation characteristics are not captured by
training data for the signal-processing models. At a 10\,ms hop, CREPE produces
16,269 frames over the 162.7-second recording, of which 9,912 (60.9\%) exceed
the confidence threshold of 0.5 and are treated as voiced. Xu~(1999) spike
removal is applied to each syllable nucleus measurement after CREPE, as a safety
net for any residual artefacts.

\subsection{Comparative tracker evaluation on Daakie}

For reference, three additional pitch trackers were applied to the same recording
and their outputs stored as parallel prosody tiers in the TextGrid. Table~%
\ref{tab:daakie_trackers} summarises the voiced frame counts and the proportion
of syllables receiving an unknown label (\texttt{?}), which indicates that the
tracker found insufficient voiced frames in the nucleus.

\begin{table}[h]
\centering
\begin{tabular}{lrrr}
\hline
Tracker & Voiced frames & Unknown syllables & Notes \\
\hline
CREPE (full, 10\,ms)  & 9{,}912  & 4 (0.9\%)  & best accuracy \\
pYIN (librosa)        & 16{,}634 & 11 (2.4\%) & reliable V/UV detection \\
Praat AC              & 16{,}784 & 2 (0.4\%)  & known octave errors \\
YIN (librosa)         & 27{,}350 & 0 (0.0\%)  & no V/UV detection --- unreliable \\
\hline
\end{tabular}
\caption{Pitch tracker comparison on the DoReCo Daakie recording (162.7\,s,
453 syllables). V/UV = voiced/unvoiced. YIN's zero unknown rate reflects the
absence of any voiced/unvoiced detector: it returns a pitch estimate for every
frame including unvoiced consonants and silence, making its high voiced-frame
count an artefact rather than a measure of accuracy.}
\label{tab:daakie_trackers}
\end{table}

% ENDED ADDING


\section{Results}

\subsection{Objective Performance Metrics}

\subsection{Perceived Performance Metrics}

\subsection{XMetrics}

\subsection{Sentiment Analysis (RAVDESS)}


%\chapter{Discussion}

%\section{Integrating Theory into Practice}

%\subsection{Evaluating SpeechPrint: Demand, Limitations, and Outcomes}

%\subsection{Independent Practice and Learner Implications}

%\subsection{SpeechPrint’s Limitations and Future Directions}

%\section{Conclusions}

%\chapter{Future Horizons: From Annotation to Speech Authoring}
%\label{ch:future}

%The work presented in this thesis demonstrates that the fundamental bottleneck in modern speech technology is not the \textbf{extraction} of signal, but its \textbf{representation}. While current Automatic Speech Recognition (ASR) systems perform a form of "lossy compression" by flattening the multidimensional human voice into a one-dimensional lexical stream \cite{krifkadobes2026}, \textbf{SpeechPrint} serves as a Universal Metadata Layer that recovers the discarded colors and textures of human intent. Looking forward, the evolution of this symbolic layer points toward a transition from passive analysis to active authoring, grounded in \textbf{vowel-centric reasoning} and \textbf{verifiable alignment}.

%\section{The Vowel-Nuclei Paradigm and Symbolic Grounding}
%A critical advancement for the next iteration of SpeechPrint is the full adoption of the \textit{Carrier Hypothesis}. Recent research in \textit{VowelPrompt} (2026) asserts that vowels serve as the primary stable carriers of affective prosody, hosting the most salient pitch and energy patterns while consonantal regions remain spectrally unstable \cite{wang2026vowelprompt}. 

%By narrowing "etic" extraction specifically to vowel nuclei, future versions of the SpeechPrint pipeline can achieve a significantly higher signal-to-noise ratio for emotion recognition and prosodic highlighting. This allows the symbolic layer (/, \textbackslash, –, *) to move beyond mere annotation and into the realm of \textbf{paralinguistic-aware Spoken Language Models (SLMs)}. Augmenting transcripts with fine-grained, vowel-level natural language descriptors enables Large Language Models to "hear" emotions from text by reasoning over pitch slope and intensity variation rather than consuming opaque, uninterpretable acoustic embeddings \cite{wang2026vowelprompt, wu2025beyond}.

%\section{The Speech DAW: Disentanglement and Active Authoring}
%The current prototype of SpeechPrint facilitates a "literacoustic" reading of text, but its ultimate potential lies in transitioning into a \textbf{Speech DAW (Digital Audio Workstation)} for authoring intent. Traditional tools separate text (ELAN) and prosody (Praat), whereas SpeechPrint investigates how these can be reintroduced through a unified representation \cite{krifkadobes2026, paschen2020doreco}.

%By adopting disentangled speech architectures—such as those pioneered by the \textit{Dis-Vector} project—speech can be decomposed into independent "latents" for content, pitch, rhythm, and timbre \cite{disvector2026}. In this future workstation, a user—whether a poet, voice actor, or educator—could click a specific word and "re-pitch" or "re-rhythm" the output using SpeechPrint's symbolic layer as the control surface. This treats speech as \textbf{Dynamic Vector Graphics} not just for scientific visualization, but for zero-shot speech editing and style transfer where phonetic units become navigable, editable objects \cite{pataca2021speech, mccurdy2016poemage}.

%\section{Verifiable Intent and Clinical Impact}
%As generative AI approaches "perfect" prosodic output, the role of a symbolic representation becomes one of \textbf{verification and reasoning} rather than just generation. Future developments should leverage Reinforcement Learning with Verifiable Reward (RLVR) to ensure that the "sonic topology" generated by an AI aligns precisely with a human researcher's goals \cite{deepseek2025r1, wang2026vowelprompt}.

%This has profound implications for \textbf{Oral Reading Fluency (ORF)} and clinical linguistics. It allows for the automated scoring of student expressiveness against established benchmarks, such as the NAEP Fluency Scale, providing immediate and interpretable feedback that standard "flat" ASR cannot offer \cite{sammit2022automated, naep2005}. By bridging the gap between raw audio and human-readable metadata, SpeechPrint provides the transparent interface required for humans to remain the primary authors of their own vocal intent.

%\section{Conclusion: Preservation as Originality}
%The "SpeechPrint" philosophy argues that \textbf{Speech $\neq$ Text}. The originality of this work lies in the refusal to accept the information loss inherent in standard transcription. By elevating prosody to a first-class annotation layer, we preserve the "phonetic fingerprint" of identity and the sonic topology of human art \cite{krifkadobes2026, clement2013sounding}. In an era of increasingly black-box AI, the symbolic grounding of prosody ensures that the "voice" of the speaker survives the transition to the digital page.



\newpage



% ADDED HERE — future directions: TTS resynthesis, prosody DAW, classification model
% (from discussion with Lonce: he mentioned TTS and prosody control; also my own idea for a Speech DAW)

\section{Future Directions}
\label{sec:future}

\subsection{Prosody Resynthesis: From Annotation to Sound}
\label{sec:resynthesis}

A natural extension of the SpeechPrint pipeline is to close the loop between
analysis and synthesis. The current system moves in one direction: from audio
to symbolic annotation. The inverse direction -- from symbolic annotation to
audio -- would constitute a \emph{prosody resynthesis} capability. This concept
is illustrated in Figure~\ref{fig:resynthesis_pipeline}.

The proposed pipeline would work as follows. First, a recording is processed by
the existing SpeechPrint pipeline to produce a TextGrid with a prosody tier.
Second, a user modifies the prosody symbols in the TextGrid, for example,
changing a falling accent (\texttt{*\textbackslash}) to a rising one
(\texttt{*/}) on the word \textit{Mary} to produce the reading ``MARY flew to
Milan yesterday'' rather than ``Mary FLEW to Milan yesterday.'' Third, the
modified symbolic annotation is translated into prosody control parameters that
are passed to a speech synthesis system. The synthesis system produces a new
audio file with the requested prosodic pattern while preserving the voice quality
and segmental content of the original.

The linguistic term for this operation is \emph{prosody transfer} or
\emph{intonation resynthesis}. It has been studied in the context of
human-to-human style transfer and as a component of expressive text-to-speech
synthesis \cite{skerry2018towards}. The key challenge is the mapping from
symbolic prosody (the / and \textbackslash\ labels of SpeechPrint's tier) to
quantitative F0 targets that a synthesis system can use. One approach is to
learn this mapping from a corpus of annotated speech; another is to use a
parametric F0 model in which the symbolic labels directly control the slope,
height, and velocity of F0 movements.

A practical implementation using currently available tools would use the
Coqui TTS framework or a similar system that exposes duration and pitch control
parameters. Coqui XTTS, for example, allows voice cloning from a short reference
segment and accepts prosodic modifications through its Python API. The SpeechPrint
symbolic tier would be compiled into a sequence of pitch targets at word or
syllable boundaries, which the synthesis system would interpolate into a smooth
F0 contour. This would allow a researcher to hear the prosodic consequences of an
annotation change without recording a new utterance -- a capability that could be
valuable both for validation of the annotation and for generating stimuli in
prosody experiments.

\begin{figure}[htbp]
\centering
\fbox{\parbox{0.9\linewidth}{\centering\vspace{2mm}
\textbf{Analysis direction (current):} \\[2mm]
WAV $\rightarrow$ Whisper $\rightarrow$ MFA $\rightarrow$ pYIN $\rightarrow$
Symbolic layer (/ \textbackslash\, * - \_) $\rightarrow$ TextGrid \\[4mm]
\textbf{Resynthesis direction (proposed):} \\[2mm]
TextGrid (edited symbols) $\rightarrow$ F0 target compiler $\rightarrow$
Coqui TTS / XTTS $\rightarrow$ Resynthesized WAV \\[2mm]
}}
\caption{The proposed bidirectional SpeechPrint pipeline. The analysis direction
is fully implemented; the resynthesis direction is a design concept for future
work. The loop would allow a user to modify prosody symbols and hear the result
without re-recording.}
\label{fig:resynthesis_pipeline}
\end{figure}

\subsection{A Speech DAW: Interactive Prosody Manipulation}
\label{sec:speech_daw}

Taking the resynthesis concept further, a \emph{speech digital audio workstation}
(Speech DAW) would present the TextGrid annotation tiers as an editable timeline,
similar in concept to how a music DAW presents MIDI notes. The user would be able
to click on any syllable label in the prosody tier and change it, drag the onset
or offset of a pitch movement, raise or lower the height symbol from \texttt{\_}
to \texttt{-}, or add or remove an accent marker \texttt{*}. Each edit would
trigger a re-synthesis of the affected region, allowing real-time auditory
feedback.

The design is motivated by a gap in current tools. Praat allows the manual
manipulation of a pitch track at the acoustic level, which requires expert
knowledge of what a changed F0 contour sounds like. The proposed Speech DAW
would operate at the symbolic level, where the units (/ \textbackslash\, *) are
interpretable by a non-specialist. This lowers the barrier for users in
conversation analysis or field linguistics who want to test hypotheses about
prosodic structure by listening to counterfactual examples.

\subsection{Automatic Prosody Classification}
\label{sec:prosody_classifier}

A third future direction is the development of a classifier that assigns one of a
small number of prosodic categories to each sentence or utterance. The motivation
is that a single global label -- for example, whether a sentence is a statement, a
question, an emphatic assertion, a continuation, or a boundary -- is often more
useful for annotation workflows than a fine-grained per-syllable sequence of
symbols. The classification would be trained on existing annotated corpora and
applied to new recordings.

A candidate set of five prosodic categories, sufficient to cover the most
frequent intonational events in English and German, might be:
\begin{enumerate}
\item \textbf{Declarative}: broad-focus statement with final fall.
\item \textbf{Interrogative}: yes/no question with final rise.
\item \textbf{Narrow focus}: prominent accent on a single content word,
      marking contrast or correction.
\item \textbf{Continuation}: non-final intonational phrase with
      non-falling boundary tone, signalling that more speech follows.
\item \textbf{Exclamation}: emphatic or expressive contour, typically with
      high onset and rapid fall.
\end{enumerate}

This classification scheme is a simplification of the GToBI inventory but
captures the categories most relevant to the use cases described in
Chapter~\ref{scenarios}. The classifier would take the sequence of SpeechPrint
symbolic labels as input features, supplemented by the mean F0, F0 range, and
final boundary tone of the utterance. The output category would be displayed as
an additional tier in the TextGrid, providing a sentence-level interpretive
summary alongside the syllable-level labels.

\chapter{Discussion and Future Directions}
\label{ch:discussion}


%\chapter{Discussion}

%\section{Integrating Theory into Practice}

%\subsection{Evaluating SpeechPrint: Demand, Limitations, and Outcomes}

%\subsection{Independent Practice and Learner Implications}

%\subsection{SpeechPrint’s Limitations and Future Directions}

%\section{Conclusions}

%\chapter{Future Horizons: From Annotation to Speech Authoring}
%\label{ch:future}

%The work presented in this thesis demonstrates that the fundamental bottleneck in modern speech technology is not the \textbf{extraction} of signal, but its \textbf{representation}. While current Automatic Speech Recognition (ASR) systems perform a form of "lossy compression" by flattening the multidimensional human voice into a one-dimensional lexical stream \cite{krifkadobes2026}, \textbf{SpeechPrint} serves as a Universal Metadata Layer that recovers the discarded colors and textures of human intent. Looking forward, the evolution of this symbolic layer points toward a transition from passive analysis to active authoring, grounded in \textbf{vowel-centric reasoning} and \textbf{verifiable alignment}.

%\section{The Vowel-Nuclei Paradigm and Symbolic Grounding}
%A critical advancement for the next iteration of SpeechPrint is the full adoption of the \textit{Carrier Hypothesis}. Recent research in \textit{VowelPrompt} (2026) asserts that vowels serve as the primary stable carriers of affective prosody, hosting the most salient pitch and energy patterns while consonantal regions remain spectrally unstable \cite{wang2026vowelprompt}. 

%By narrowing "etic" extraction specifically to vowel nuclei, future versions of the SpeechPrint pipeline can achieve a significantly higher signal-to-noise ratio for emotion recognition and prosodic highlighting. This allows the symbolic layer (/, \textbackslash, –, *) to move beyond mere annotation and into the realm of \textbf{paralinguistic-aware Spoken Language Models (SLMs)}. Augmenting transcripts with fine-grained, vowel-level natural language descriptors enables Large Language Models to "hear" emotions from text by reasoning over pitch slope and intensity variation rather than consuming opaque, uninterpretable acoustic embeddings \cite{wang2026vowelprompt, wu2025beyond}.

%\section{The Speech DAW: Disentanglement and Active Authoring}
%The current prototype of SpeechPrint facilitates a "literacoustic" reading of text, but its ultimate potential lies in transitioning into a \textbf{Speech DAW (Digital Audio Workstation)} for authoring intent. Traditional tools separate text (ELAN) and prosody (Praat), whereas SpeechPrint investigates how these can be reintroduced through a unified representation \cite{krifkadobes2026, paschen2020doreco}.

%By adopting disentangled speech architectures—such as those pioneered by the \textit{Dis-Vector} project—speech can be decomposed into independent "latents" for content, pitch, rhythm, and timbre \cite{disvector2026}. In this future workstation, a user—whether a poet, voice actor, or educator—could click a specific word and "re-pitch" or "re-rhythm" the output using SpeechPrint's symbolic layer as the control surface. This treats speech as \textbf{Dynamic Vector Graphics} not just for scientific visualization, but for zero-shot speech editing and style transfer where phonetic units become navigable, editable objects \cite{pataca2021speech, mccurdy2016poemage}.

%\section{Verifiable Intent and Clinical Impact}
%As generative AI approaches "perfect" prosodic output, the role of a symbolic representation becomes one of \textbf{verification and reasoning} rather than just generation. Future developments should leverage Reinforcement Learning with Verifiable Reward (RLVR) to ensure that the "sonic topology" generated by an AI aligns precisely with a human researcher's goals \cite{deepseek2025r1, wang2026vowelprompt}.

%This has profound implications for \textbf{Oral Reading Fluency (ORF)} and clinical linguistics. It allows for the automated scoring of student expressiveness against established benchmarks, such as the NAEP Fluency Scale, providing immediate and interpretable feedback that standard "flat" ASR cannot offer \cite{sammit2022automated, naep2005}. By bridging the gap between raw audio and human-readable metadata, SpeechPrint provides the transparent interface required for humans to remain the primary authors of their own vocal intent.

%\section{Conclusion: Preservation as Originality}
%The "SpeechPrint" philosophy argues that \textbf{Speech $\neq$ Text}. The originality of this work lies in the refusal to accept the information loss inherent in standard transcription. By elevating prosody to a first-class annotation layer, we preserve the "phonetic fingerprint" of identity and the sonic topology of human art \cite{krifkadobes2026, clement2013sounding}. In an era of increasingly black-box AI, the symbolic grounding of prosody ensures that the "voice" of the speaker survives the transition to the digital page.



\newpage



% ADDED HERE — future directions: TTS resynthesis, prosody DAW, classification model
% (from discussion with Lonce: he mentioned TTS and prosody control; also my own idea for a Speech DAW)

\section{Future Directions}
\label{sec:future}

\subsection{Prosody Resynthesis: From Annotation to Sound}
\label{sec:resynthesis}

A natural extension of the SpeechPrint pipeline is to close the loop between
analysis and synthesis. The current system moves in one direction: from audio
to symbolic annotation. The inverse direction -- from symbolic annotation to
audio -- would constitute a \emph{prosody resynthesis} capability. This concept
is illustrated in Figure~\ref{fig:resynthesis_pipeline}.

The proposed pipeline would work as follows. First, a recording is processed by
the existing SpeechPrint pipeline to produce a TextGrid with a prosody tier.
Second, a user modifies the prosody symbols in the TextGrid, for example,
changing a falling accent (\texttt{*\textbackslash}) to a rising one
(\texttt{*/}) on the word \textit{Mary} to produce the reading ``MARY flew to
Milan yesterday'' rather than ``Mary FLEW to Milan yesterday.'' Third, the
modified symbolic annotation is translated into prosody control parameters that
are passed to a speech synthesis system. The synthesis system produces a new
audio file with the requested prosodic pattern while preserving the voice quality
and segmental content of the original.

The linguistic term for this operation is \emph{prosody transfer} or
\emph{intonation resynthesis}. It has been studied in the context of
human-to-human style transfer and as a component of expressive text-to-speech
synthesis \cite{skerry2018towards}. The key challenge is the mapping from
symbolic prosody (the / and \textbackslash\ labels of SpeechPrint's tier) to
quantitative F0 targets that a synthesis system can use. One approach is to
learn this mapping from a corpus of annotated speech; another is to use a
parametric F0 model in which the symbolic labels directly control the slope,
height, and velocity of F0 movements.

A practical implementation using currently available tools would use the
Coqui TTS framework or a similar system that exposes duration and pitch control
parameters. Coqui XTTS, for example, allows voice cloning from a short reference
segment and accepts prosodic modifications through its Python API. The SpeechPrint
symbolic tier would be compiled into a sequence of pitch targets at word or
syllable boundaries, which the synthesis system would interpolate into a smooth
F0 contour. This would allow a researcher to hear the prosodic consequences of an
annotation change without recording a new utterance -- a capability that could be
valuable both for validation of the annotation and for generating stimuli in
prosody experiments.

\begin{figure}[htbp]
\centering
\fbox{\parbox{0.9\linewidth}{\centering\vspace{2mm}
\textbf{Analysis direction (current):} \\[2mm]
WAV $\rightarrow$ WhisperX $\rightarrow$ MFA / hardcoded IPA
$\rightarrow$ \{pYIN (read speech) $|$ CREPE (archival)\}
$\rightarrow$ Symbolic layer (/ \textbackslash\, * {\={}} \_)
$\rightarrow$ TextGrid \\[4mm]
\textbf{Resynthesis direction (proposed):} \\[2mm]
TextGrid (edited symbols) $\rightarrow$ F0 target compiler $\rightarrow$
Coqui TTS / XTTS $\rightarrow$ Resynthesized WAV \\[2mm]
}}
\caption{The proposed bidirectional SpeechPrint pipeline. The analysis direction
is fully implemented; the resynthesis direction is a design concept for future
work. The loop would allow a user to modify prosody symbols and hear the result
without re-recording.}
\label{fig:resynthesis_pipeline}
\end{figure}

\subsection{A Speech DAW: Interactive Prosody Manipulation}
\label{sec:speech_daw}

Taking the resynthesis concept further, a \emph{speech digital audio workstation}
(Speech DAW) would present the TextGrid annotation tiers as an editable timeline,
similar in concept to how a music DAW presents MIDI notes. The user would be able
to click on any syllable label in the prosody tier and change it, drag the onset
or offset of a pitch movement, raise or lower the height symbol from \texttt{\_}
to \texttt{{\={}}}, or add or remove an accent marker \texttt{*}. Each edit would
trigger a re-synthesis of the affected region, allowing real-time auditory
feedback.

The design is motivated by a gap in current tools. Praat allows the manual
manipulation of a pitch track at the acoustic level, which requires expert
knowledge of what a changed F0 contour sounds like. The proposed Speech DAW
would operate at the symbolic level, where the units (/ \textbackslash\, *) are
interpretable by a non-specialist. This lowers the barrier for users in
conversation analysis or field linguistics who want to test hypotheses about
prosodic structure by listening to counterfactual examples.

\subsection{Automatic Prosody Classification}
\label{sec:prosody_classifier}

A third future direction is the development of a classifier that assigns one of a
small number of prosodic categories to each sentence or utterance. The motivation
is that a single global label -- for example, whether a sentence is a statement, a
question, an emphatic assertion, a continuation, or a boundary -- is often more
useful for annotation workflows than a fine-grained per-syllable sequence of
symbols. The classification would be trained on existing annotated corpora and
applied to new recordings.

A candidate set of five prosodic categories, sufficient to cover the most
frequent intonational events in English and German, might be:
\begin{enumerate}
\item \textbf{Declarative}: broad-focus statement with final fall.
\item \textbf{Interrogative}: yes/no question with final rise.
\item \textbf{Narrow focus}: prominent accent on a single content word,
      marking contrast or correction.
\item \textbf{Continuation}: non-final intonational phrase with
      non-falling boundary tone, signalling that more speech follows.
\item \textbf{Exclamation}: emphatic or expressive contour, typically with
      high onset and rapid fall.
\end{enumerate}

This classification scheme is a simplification of the GToBI inventory but
captures the categories most relevant to the use cases described in
Chapter~\ref{scenarios}. The classifier would take the sequence of SpeechPrint
symbolic labels as input features, supplemented by the mean F0, F0 range, and
final boundary tone of the utterance. The output category would be displayed as
an additional tier in the TextGrid, providing a sentence-level interpretive
summary alongside the syllable-level labels.

\subsection{Bias, Projection, and Cross-Lingual Phonological Approximation}
\label{sec:bias_projection}

A recurring observation in this thesis is that applying a speech model trained
on one language to a recording in an unrelated language does not produce random
noise --- it produces something coherent but wrong. When Whisper's Italian
acoustic model was applied to the Daakie recording, it returned phonetically
plausible Italian word sequences, not silence or unintelligible garbage
(Appendix~\ref{app:daakie}). The model could not understand Daakie, but it
projected its learned phonological structure onto the signal and produced the
nearest thing in its training distribution. This is model bias operating as a
generative mechanism: the failure is systematic and, viewed from a different
angle, compositionally interesting. Each reinterpretation layer introduces a
transformation that is determined by what the model knows and what it does not,
producing emergent changes in apparent timbre, prosody, and intelligibility that
are not random but structured by the model's prior. The SpeechPrint prosody
layer applied to such a mistranscribed signal assigns labels that are acoustically
correct but linguistically misaligned, creating a kind of prosodic ghost of the
original utterance. This phenomenon --- the gap between what a model hears and
what was said --- represents a potential compositional space for work at the
intersection of speech technology and sound art, and connects to a broader
research thread on how bias and error in computational listening shape the
transformation of voice across systems.

% ENDED ADDING

\newpage




